\documentclass[aspectratio=169,11pt]{beamer}
% Shared Beamer style for practice contest solution walkthroughs.
\usepackage[utf8]{inputenc}
\usepackage[T1]{fontenc}
\usepackage{lmodern}
\usepackage{amsmath,amssymb}
\usepackage{xcolor}
\usepackage{hyperref}
\usepackage{listings}

\definecolor{CPNavy}{HTML}{172033}
\definecolor{CPBlue}{HTML}{2563EB}
\definecolor{CPGreen}{HTML}{16A34A}
\definecolor{CPOrange}{HTML}{F97316}
\definecolor{CPGray}{HTML}{64748B}
\definecolor{CPLight}{HTML}{F8FAFC}
\definecolor{CPCodeBg}{HTML}{F1F5F9}
\definecolor{CPCodeRule}{HTML}{CBD5E1}

\usetheme{default}
\usefonttheme{professionalfonts}
\setbeamertemplate{navigation symbols}{}
\setbeamersize{text margin left=0.65cm,text margin right=0.65cm}
\setbeamertemplate{itemize item}{\color{CPBlue}$\blacktriangleright$}
\setbeamertemplate{itemize subitem}{\color{CPGray}$\triangleright$}

\setbeamercolor{normal text}{fg=CPNavy,bg=white}
\setbeamercolor{frametitle}{fg=white,bg=CPNavy}
\setbeamercolor{title}{fg=white,bg=CPNavy}
\setbeamercolor{subtitle}{fg=CPLight,bg=CPNavy}
\hypersetup{colorlinks=true,urlcolor=CPBlue}

\setbeamertemplate{frametitle}{%
  \nointerlineskip%
  \begin{beamercolorbox}[wd=\paperwidth,ht=0.88cm,dp=0.22cm,leftskip=0.55cm,rightskip=0.35cm]{frametitle}
    \usebeamerfont{frametitle}\insertframetitle\hfill\small\insertframenumber/\inserttotalframenumber
  \end{beamercolorbox}%
}

\setbeamertemplate{footline}{%
  \leavevmode%
  \hbox{%
    \begin{beamercolorbox}[wd=.58\paperwidth,ht=2.4ex,dp=1ex,leftskip=.5cm]{author in head/foot}%
      \color{CPGray}\insertshorttitle
    \end{beamercolorbox}%
    \begin{beamercolorbox}[wd=.42\paperwidth,ht=2.4ex,dp=1ex,rightskip=.5cm plus1fil]{date in head/foot}%
      \hfill\color{CPGray}\insertshortauthor
    \end{beamercolorbox}%
  }%
}

\setbeamertemplate{title page}{%
  \vspace*{1.25cm}
  \begin{beamercolorbox}[wd=\paperwidth,sep=0.65cm,leftskip=0.15cm,rightskip=0.15cm]{title}
    {\color{white}\Huge\bfseries\inserttitle\par}
    \vspace{0.35cm}
    {\color{CPLight}\Large\insertsubtitle\par}
    \vspace{0.6cm}
    {\color{CPLight}\large\insertauthor\par}
  \end{beamercolorbox}
  \vfill
}

\newcommand{\sectiontag}[1]{\textbf{\color{CPBlue}#1}}
\newenvironment{tightitemize}{\begin{itemize}\small\setlength{\itemsep}{0.18cm}}{\end{itemize}}
\lstdefinestyle{cpcode}{
  language=Python,
  basicstyle=\ttfamily\tiny,
  keywordstyle=\color{CPBlue}\bfseries,
  commentstyle=\color{CPGray}\itshape,
  stringstyle=\color{CPGreen},
  backgroundcolor=\color{CPCodeBg},
  frame=single,
  framerule=0.25pt,
  rulecolor=\color{CPCodeRule},
  showstringspaces=false,
  columns=fullflexible,
  keepspaces=true,
  breaklines=true,
  breakatwhitespace=false,
  tabsize=4,
  xleftmargin=0.08cm,
  xrightmargin=0.08cm,
  aboveskip=0.08cm,
  belowskip=0.08cm
}


\title[hurricanedanger]{Hurricane Danger!}
\subtitle{Day 5 solution walkthrough}
\author[Solution Deck]{Competitive Programming Bootcamp}
\date{}

\begin{document}

\begin{frame}[plain]
  \titlepage
\end{frame}

% BEGIN GENERATED STATEMENT INTERPRETATION
\begin{frame}{Interpret the Word Problem}
\small
\sectiontag{Statement excerpts:}
\begin{tightitemize}
  \item \textit{``city that is closest to the path''}
  \item \textit{``travels in a straight line''}
\end{tightitemize}

\vspace{0.18cm}
\sectiontag{Programming interpretation:}
\begin{tightitemize}
  \item The hurricane path is an infinite line through two observed points.
  \item Every city's risk is determined by perpendicular distance to that line.
  \item Since the denominator is shared, compare only the absolute line-equation numerator.
\end{tightitemize}
\end{frame}
% END GENERATED STATEMENT INTERPRETATION







\begin{frame}{Problem Model}
\small
\sectiontag{Textbook connection:} CP4 7.2.2 Lines

\vspace{0.25cm}
\sectiontag{Core technique:} Point-to-line distance comparison

\vspace{0.25cm}
\begin{tightitemize}
  \item The hurricane path is an infinite line through two observed points.
  \item Each city has a name and coordinates.
  \item Print all cities closest to the hurricane path.
\end{tightitemize}
\end{frame}

\begin{frame}{How to Recognize the Approach}
\begin{tightitemize}
  \item Distance is from a point to an infinite line, not to a segment.
  \item The denominator of the distance formula is the same for every city in a test case.
  \item Comparing only the numerator avoids floating-point equality problems.
\end{tightitemize}
\end{frame}

\begin{frame}{Algorithm}
\begin{tightitemize}
  \item Convert the two hurricane points into A*x + B*y + C = 0.
  \item For each city, compute abs(A*x + B*y + C).
  \item If this value is smaller than the best, reset the answer list.
  \item If it ties the best, append the city name.
  \item Print tied names in input order.
\end{tightitemize}
\end{frame}

\begin{frame}{Walkthrough of the Key State}
\begin{tightitemize}
  \item A = y2 - y1, B = -(x2 - x1), C = x2*y1 - y2*x1.
  \item The full distance divides by sqrt(A\textasciicircum{}2 + B\textasciicircum{}2).
  \item Since that denominator is constant, numerator comparison gives the same ranking.
\end{tightitemize}
\end{frame}

\begin{frame}{Why the Algorithm Is Correct}
\begin{tightitemize}
  \item The standard line equation represents exactly the hurricane path.
  \item Point-to-line distance ranks cities by abs(Ax + By + C) when the line is fixed.
  \item The algorithm keeps exactly all cities with the smallest such value.
\end{tightitemize}
\end{frame}

\begin{frame}{Complexity and Implementation Notes}
\small
\sectiontag{Complexity:} O(M) time per test case and O(M) memory in the all-tied case.

\vspace{0.3cm}
\begin{tightitemize}
  \item Using integers avoids precision issues when checking ties.
  \item Keep city names in input order when distances tie.
\end{tightitemize}
\end{frame}



% BEGIN GENERATED CODE WALKTHROUGH
\begin{frame}[fragile]{Code: Build the Line Equation}
\scriptsize
\sectiontag{Source lines:} \texttt{hurricanedanger.py:45--65}

\vspace{0.08cm}
\begin{columns}[T,totalwidth=\textwidth]
\begin{column}{0.58\textwidth}
\begin{lstlisting}[style=cpcode]
import sys

def main():
    input = sys.stdin.buffer.readline

    test_cases = int(input())

    for _ in range(test_cases):
        x1, y1, x2, y2 = map(int, input().split())

        cities = int(input())

        A = y2 - y1
        B = -(x2 - x1)
        C = x2 * y1 - y2 * x1
\end{lstlisting}
\end{column}
\begin{column}{0.38\textwidth}
\sectiontag{What this chunk does}
\begin{tightitemize}
  \item Each case begins with two points on the hurricane path.
  \item The line is stored as A*x + B*y + C = 0.
  \item Those coefficients are reused for every city in the case.
\end{tightitemize}
\end{column}
\end{columns}
\end{frame}

\begin{frame}[fragile]{Code: Measure City Distance}
\scriptsize
\sectiontag{Source lines:} \texttt{hurricanedanger.py:67--83}

\vspace{0.08cm}
\begin{columns}[T,totalwidth=\textwidth]
\begin{column}{0.58\textwidth}
\begin{lstlisting}[style=cpcode]
best_distance = None
answer = []

for _ in range(cities):
    data = input().split()

    name = data[0].decode()
    x = int(data[1])
    y = int(data[2])

    distance = abs(A * x + B * y + C)
\end{lstlisting}
\end{column}
\begin{column}{0.38\textwidth}
\sectiontag{What this chunk does}
\begin{tightitemize}
  \item Each city gives a name and integer coordinates.
  \item The distance variable is the absolute numerator of point-to-line distance.
  \item The omitted denominator is identical for all cities in the test case.
\end{tightitemize}
\end{column}
\end{columns}
\end{frame}

\begin{frame}[fragile]{Code: Keep Closest Ties}
\scriptsize
\sectiontag{Source lines:} \texttt{hurricanedanger.py:85--98}

\vspace{0.08cm}
\begin{columns}[T,totalwidth=\textwidth]
\begin{column}{0.58\textwidth}
\begin{lstlisting}[style=cpcode]
            if best_distance is None or distance < best_distance:
                best_distance = distance
                answer = [name]

            elif distance == best_distance:
                answer.append(name)

        print(" ".join(answer))

main()
\end{lstlisting}
\end{column}
\begin{column}{0.38\textwidth}
\sectiontag{What this chunk does}
\begin{tightitemize}
  \item A smaller distance replaces the current answer list.
  \item An equal distance appends the city name in input order.
  \item The final line prints all most-dangerous cities separated by spaces.
\end{tightitemize}
\end{column}
\end{columns}
\end{frame}
% END GENERATED CODE WALKTHROUGH

\begin{frame}{Source}
\small
\sectiontag{Solution file:} \texttt{practice\_contests/day\_5/hurricanedanger.py}

\vspace{0.3cm}
\sectiontag{Problem:} \url{https://open.kattis.com/problems/hurricanedanger}

\vspace{0.45cm}
Use this deck with the source code open beside it. The slides explain the reasoning; the code shows the exact input handling and output formatting.
\end{frame}

\end{document}
